%%%%%%%%%%%%%%%%%%%%%%%%%%%%%%%%%%%%%%%%%%%%%%%%%%%%%%%%%%%%%%%
% SECTION 3.5: THEORETICAL PROPERTIES OF THE O/R INDEX
% Insert this after Section 3.4 (Theoretical Motivation)
%%%%%%%%%%%%%%%%%%%%%%%%%%%%%%%%%%%%%%%%%%%%%%%%%%%%%%%%%%%%%%%

\subsection{Theoretical Properties}
\label{subsec:theory_properties}

While our primary contribution is empirical validation, we establish formal properties that characterize when and why the O/R Index succeeds as a coordination predictor.

\begin{proposition}[Range and Extremes]
\label{prop:range_extremes}
Let $\pi$ be a policy with discrete observation space $\mathcal{O}$ and finite discrete action space $\mathcal{A}$. Then the O/R Index satisfies:
\begin{enumerate}
    \item $\mathrm{OR}(\pi) \in [-1, \infty)$.
    \item $\mathrm{OR}(\pi) = -1$ whenever $\pi$ is deterministic given observations (i.e., $\forall o \in \mathcal{O}$, $\pi(a|o)$ is a point mass).
    \item $\mathrm{OR}(\pi) = 0$ if and only if actions are conditionally independent of observations in the ANOVA sense: $\Var(P(a|o)) = \Var(P(a))$.
    \item $\mathrm{OR}(\pi) > 0$ can occur when observation bins have highly unequal sizes and within-bin variance exceeds total variance due to unweighted bin averaging.
\end{enumerate}
\end{proposition}

\begin{proof}[Proof sketch]
See Appendix~\ref{app:proofs_discrete} for complete proofs. The range follows from: (1) deterministic policies yield zero within-bin variance, giving $\mathrm{OR} = 0/\Var(P(a)) - 1 = -1$; (2) when observations provide no information, within-bin variance equals total variance, giving $\mathrm{OR} = 0$; (3) unweighted averaging over bins with unequal sample sizes can produce $\Var(P(a|o)) > \Var(P(a))$, yielding $\mathrm{OR} > 0$.
\end{proof}

\textbf{Remark on Range.} In general, the O/R Index satisfies $\mathrm{OR}(\pi) \in [-1, \infty)$. The lower bound $-1$ is attained by deterministic observation-based policies. The neutral point $0$ corresponds to policies whose actions are independent of observations. Values greater than $0$ arise when conditioning on observations increases the apparent variability of the action distribution, typically due to skewed observation bins or undersampled states. However, many natural policy constructions (such as noise-mixing in Proposition~\ref{prop:monotonicity_noise}) produce O/R values confined to $[-1, 0]$.

\begin{proposition}[Monotonicity under Noise Mixing]
\label{prop:monotonicity_noise}
Let $\pi_0$ be a deterministic observation-based policy mapping each observation $o$ to a fixed action $a^*(o)$, and let $\pi_{\mathrm{rand}}$ be a uniform random policy over $\mathcal{A}$. Define the mixture policy:
\begin{equation}
\pi_\alpha(a|o) = (1-\alpha) \cdot \pi_0(a|o) + \alpha \cdot \pi_{\mathrm{rand}}(a)
\end{equation}
for $\alpha \in [0, 1]$. Then $\mathrm{OR}(\pi_\alpha)$ is strictly increasing in $\alpha$ for $\alpha \in (0, 1)$.
\end{proposition}

\begin{proof}[Proof sketch]
As $\alpha$ increases, we inject more private randomness (action noise independent of observations). This increases within-bin variance $\Var(P(a|o))$ while total variance $\Var(P(a))$ remains bounded. The normalized ratio increases monotonically. See Appendix~\ref{app:proofs_discrete} for the detailed calculation.
\end{proof}

\textbf{Interpretation.} Proposition~\ref{prop:range_extremes} formalizes the intuition that O/R measures behavioral consistency: deterministic policies achieve the minimum ($-1$), while policies with high conditional variance (unpredictable responses) have high O/R. Proposition~\ref{prop:monotonicity_noise} shows that adding private randomness—behavior invisible to teammates—monotonically increases O/R, formalizing why erratic behavior hinders coordination.

\textbf{Toy Example: 2×2 Matrix Game.} Consider a coordination game with two states (observations) $o_1, o_2$ and two actions $\{L, R\}$:
\begin{center}
\begin{tabular}{c|cc}
& $a_1 = L$ & $a_1 = R$ \\
\hline
$a_2 = L$ & 1, 1 & 0, 0 \\
$a_2 = R$ & 0, 0 & 1, 1 \\
\end{tabular}
\end{center}

\textbf{Policy A (Deterministic):} Agent plays $L$ when $o = o_1$ and $R$ when $o = o_2$. Within each observation bin, variance is zero. Thus $\Var(P(a|o)) = 0$ and $\mathrm{OR} = -1$. This is maximally predictable.

\textbf{Policy B (Partially Noisy):} Mix Policy A with uniform random noise at rate $\alpha = 0.5$. Empirical validation (see Appendix~\ref{app:proofs_discrete}) shows this yields $\mathrm{OR} \approx -0.25$, confirming Proposition~\ref{prop:monotonicity_noise}.

\textbf{Policy C (Fully Random):} Uniform random policy over $\{L, R\}$ independent of observations. This yields $\mathrm{OR} \approx 0$, as actions provide no signal about observations.

This toy example validates Proposition~\ref{prop:monotonicity_noise}: in a 2×2 coordination game, our noise-mixing construction produces O/R values confined to $[-1, 0]$, smoothly increasing from $-1.0$ (deterministic) to $\approx 0.0$ (random) as noise level increases. The absence of values $> 0$ here reflects that uniform noise mixing cannot make conditional variance exceed marginal variance in this construction; exploring the $\mathrm{OR} > 0$ regime would require different policy structures (e.g., skewed binning or highly non-uniform observation distributions).

These propositions provide theoretical grounding for our empirical finding that lower O/R predicts coordination success: teams with more deterministic, observation-consistent policies (low O/R) achieve higher coordination because teammates can reliably anticipate behavior under partial observability.

%%%%%%%%%%%%%%%%%%%%%%%%%%%%%%%%%%%%%%%%%%%%%%%%%%%%%%%%%%%%%%%
% END OF SECTION 3.5
%%%%%%%%%%%%%%%%%%%%%%%%%%%%%%%%%%%%%%%%%%%%%%%%%%%%%%%%%%%%%%%
